\section{Empirical demonstration: instrumenting the Khan (2026) volatility panel}\label{sec:demonstration}

This section reports the empirical demonstration of \texttt{mr-audit} on a real ML pipeline. We retrospectively instrument the seven-model volatility-forecasting panel of \citet{KhanVol2026}: three GARCH variants (GARCH, EGARCH, GJR-GARCH), HAR-RV, LightGBM, XGBoost, and an equal-weight ensemble, predicting 5-day-ahead realised volatility on the S\&P 500 over 980 trading days from 2022-01-03 to 2025-11-26.

\subsection{Instrumentation setup}\label{sec:demonstration_setup}

The instrumentation is retrospective. The seven-model panel of \citet{KhanVol2026} ships its forecast outputs as \texttt{forecasts\_5d.parquet} (980 rows $\times$ 7 model columns + 1 actual column) at \href{https://github.com/ayk5511/volatility-forecasting}{ayk5511/volatility-forecasting}. We do not re-run the panel's training or forecasting code. Instead, we treat each (date, model) cell of the parquet as one prediction event and write one \texttt{mr-audit} record per cell, capturing what an integrated-from-the-start instrumentation would have recorded.

For each prediction event, the recorded inputs are the five lagged realised-volatility values $(\text{rv}_{t-1}, \text{rv}_{t-2}, \text{rv}_{t-3}, \text{rv}_{t-4}, \text{rv}_{t-5})$, which are the most natural input for a 5-day-ahead vol forecast and which provide a meaningful basis for drift analysis over the test period. The \texttt{prediction.value} is the forecast value the panel actually produced for that (date, model) cell. The \texttt{model.name} is one of the seven model names; \texttt{model.version} is set to \texttt{paper2-v1} for all records to match the source release.

The instrumentation is a one-line decorator addition (see \texttt{studies/01\_instrument\_paper2.py}):

\begin{lstlisting}[language=Python]
@auditable(model_name=model_name, model_version="paper2-v1",
           horizon_days=5, feature_extractor=feature_extractor)
def predict_volatility(*, lagged_rvs, model_params=None):
    return float(model_params.get("_forecast_value", 0.0))

with AuditContext(log_path):
    for i in range(LOOKBACK_DAYS, n):
        for model_name in MODELS:
            PREDICTORS[model_name](lagged_rvs=lagged,
                                    model_params={"_forecast_value": forecast},
                                    _data_source="ayk5511/volatility-forecasting:forecasts_5d.parquet",
                                    _predicted_for_date=date.isoformat()[:10])
\end{lstlisting}

The run produces a single SQLite log file at \texttt{studies/audit\_logs/paper2\_demo.db}.

\begin{table}[H]
\centering
\caption{Instrumentation throughput on the Khan (2026) panel ($n=6825$ records). Hardware: M5 Max, 128~GB RAM. The ``bytes per record'' figure includes the SQLite indexing overhead.}\label{tab:instrumentation}
\footnotesize
\begin{tabular}{lr}
\toprule
Metric & Value \\
\midrule
Trading days in panel & 980 \\
Models per day & 7 \\
Records written ($980 - 5$ lookback) $\times 7$ & 6{,}825 \\
Wall-clock time & 302 s \\
Write throughput & 22.6 records/s \\
Log size on disk (SQLite) & 15.1 MB \\
Bytes per record (incl.\ indexes) & 2{,}266 \\
Hash-chain integrity (post-write \texttt{verify\_chain}) & valid \\
\bottomrule
\end{tabular}
\end{table}

The throughput figure (22.6 records/s) is bottlenecked by the SQLite single-writer model with synchronous writes; the package can be reconfigured for higher-throughput batch writes if needed. For a real bank-scale workload of tens of thousands of predictions per day, the current configuration is comfortable; for million-prediction-per-day pipelines, the SQLite backend would need to be replaced with a write-optimised backend (Parquet, append-only file shards, or a write-heavy database).

\subsection{Use case 1: Reproducibility audit}\label{sec:demonstration_replay}

A model risk validation analyst, given the audit log alone, must be able to reproduce any prediction the model made. We sampled 5 records uniformly at random from the 6,825-record log and replayed each via \texttt{mr\_audit.replay.replay()}, supplying the recorded inputs to the unwrapped prediction function and comparing the result to the recorded prediction.

\begin{table}[H]
\centering
\caption{Reproducibility audit results: 5 randomly-sampled records from the 6,825-record log. Tolerance for prediction equality is $10^{-12}$.}\label{tab:replay}
\footnotesize
\begin{tabular}{lllrcc}
\toprule
Record & Model & Predicted-for date & Logged prediction & Input hash match & Prediction match \\
\midrule
\texttt{[975]}  & GJR-GARCH & 2022-08-01 & 0.151071 & \checkmark & \checkmark \\
\texttt{[2617]} & Ensemble  & 2023-07-07 & 0.117039 & \checkmark & \checkmark \\
\texttt{[4116]} & GARCH     & 2024-05-14 & 0.117928 & \checkmark & \checkmark \\
\texttt{[4192]} & Ensemble  & 2024-05-29 & 0.096347 & \checkmark & \checkmark \\
\texttt{[5301]} & GJR-GARCH & 2025-01-16 & 0.159171 & \checkmark & \checkmark \\
\midrule
\multicolumn{4}{l}{\textit{Pass rate}} & \multicolumn{2}{r}{5/5 (100\%)} \\
\bottomrule
\end{tabular}
\end{table}

All 5 sampled replays match bit-identically. This is the expected behaviour for a deterministic prediction function with fully-recorded inputs: any failure here would indicate either a hash-chain breach or a code-state change between record time and replay time.

To confirm that the replay engine actually detects mismatches, we additionally verified (in \texttt{tests/test\_replay.py}) that supplying wrong inputs (e.g.\ \texttt{x=99.0} when the record was made with \texttt{x=2.0}) yields \texttt{input\_hash\_match = False} with an explanatory note in the result, and that supplying the correct inputs but a function with different logic (e.g.\ \texttt{return x * 3} when the recorded function was \texttt{return x * 2}) yields \texttt{prediction\_match = False}.

\subsection{Use case 2: Drift detection}\label{sec:demonstration_drift}

A model risk team is required by SR 11-7 §V (Ongoing Monitoring) to detect distributional shift in input data over time. We bucket the 6,825 records by the predicted-for month (47 calendar months from 2022-01 to 2025-11), use the first three months (2022 Q1, $n=399$ records) as the baseline, and apply \texttt{detect\_drift} to each subsequent month against this baseline.

We flag a month as showing drift if any of the five lagged-rv features registers Kolmogorov-Smirnov $p < 0.01$ against the baseline, or if any feature's PSI exceeds 0.25 (the credit-risk ``high'' threshold). The result is a monthly drift table.

\begin{table}[H]
\centering
\caption{Monthly drift detection results: 44 evaluated months (2022-04 through 2025-11) versus the 2022 Q1 baseline ($n=399$). The KS test uses the maximum statistic across all 5 lagged-rv features per month; the minimum p-value is shown. Both 2022 (high-vol regime) and 2023+ (lower-vol regime) months register drift relative to early 2022, consistent with the documented regime shift in Khan (2026). Excerpt of 9 of 44 months for compactness; full table at \texttt{studies/results/use\_cases\_summary.json}.}\label{tab:drift}
\footnotesize
\begin{tabular}{llrrrr}
\toprule
Month & $n$ & Max KS statistic & Min KS $p$-value & Max PSI & Drift detected \\
\midrule
2022-04 & 140 & 0.4719 & $<10^{-6}$ & high & yes \\
2022-05 & 147 & 0.6667 & $<10^{-6}$ & high & yes \\
2022-06 & 147 & 0.6266 & $<10^{-6}$ & high & yes \\
2022-09 & 147 & 0.4110 & $<10^{-6}$ & high & yes \\
\midrule
2023-07 & 147 & --     & $<10^{-6}$ & high & yes \\
2024-01 & 154 & --     & $<10^{-6}$ & high & yes \\
2024-12 & 147 & --     & $<10^{-6}$ & high & yes \\
\midrule
2025-09 & 147 & 1.0000 & $<10^{-6}$ & high & yes \\
2025-10 & 161 & 0.7651 & $<10^{-6}$ & high & yes \\
2025-11 & 126 & 0.7281 & $<10^{-6}$ & high & yes \\
\midrule
\multicolumn{5}{l}{\textit{Months with drift}}                & 44 / 44 \\
\bottomrule
\end{tabular}
\end{table}

The headline result is that drift is flagged in every month evaluated. This deserves interpretation. Each input is a lagged realised-volatility window, and realised volatility itself shifts substantially across regimes (e.g.\ Khan 2026's documented contrast between high-vol 2022 and lower-vol 2023+). A month-window of 5 lagged-RV features reflects the underlying volatility regime of that month, not a stable generative process. The KS test correctly identifies these regime-driven shifts as distributional differences. The implication for a model risk team is not that the underlying GARCH/HAR/ML models are broken; it is that the input distribution has shifted, which warrants re-validation of the models' continued fitness for purpose. This is exactly the SR 11-7 §V monitoring trigger.

A more sophisticated drift framework would also monitor \emph{prediction} drift (output distribution shift) and would distinguish covariate shift from concept drift; the package's drift module provides the input side of this analysis. The output side is a planned addition for v0.1.0 of the package.

\subsection{Use case 3: Model risk validation report}\label{sec:demonstration_validation}

A model risk validator working from the audit log alone (no external data, no live system access) should be able to produce a per-model summary report. We compute per-model statistics directly from the 6,825-record log.

\begin{table}[H]
\centering
\caption{Per-model summary report generated from the audit log alone. All seven models in the Khan (2026) panel are correctly identified, with $n=975$ predictions each (980 trading days minus 5 lookback days). Prediction means and standard deviations are summary statistics over the audit log's recorded forecast values; they are not the absolute volatility levels themselves but the model's predicted volatility values across the test period.}\label{tab:validation}
\footnotesize
\begin{tabular}{lrrr}
\toprule
Model & $n$ predictions & Prediction mean & Prediction std \\
\midrule
EGARCH    & 975 & 0.1675 & 0.0608 \\
Ensemble  & 975 & 0.1566 & 0.0657 \\
GARCH     & 975 & 0.1651 & 0.0678 \\
GJR-GARCH & 975 & 0.1611 & 0.0727 \\
HAR-RV    & 975 & 0.1540 & 0.0648 \\
LightGBM  & 975 & 0.1506 & 0.0749 \\
XGBoost   & 975 & 0.1509 & 0.0731 \\
\bottomrule
\end{tabular}
\end{table}

The library-version provenance for each model's latest record (numpy, pandas, scipy, etc.) is recorded but omitted from the table for compactness; it is in \texttt{studies/results/use\_cases\_summary.json}.

\subsection{Use case 4: Regulator export bundle}\label{sec:demonstration_export}

The fourth use case is the production of a self-contained, integrity-tied bundle that a regulator or external auditor can verify without running any \texttt{mr-audit} code. We produced the bundle from the 6,825-record log using \texttt{mr-audit export}.

\begin{table}[H]
\centering
\caption{Regulator-export bundle for the 6,825-record log. The bundle is a zip archive with five files; verification produces all-True checks across manifest hash, per-file hashes, and chain integrity.}\label{tab:bundle}
\footnotesize
\begin{tabular}{ll}
\toprule
Bundle path & \texttt{studies/audit\_logs/regulator\_bundle.zip} \\
Bundle size & 841 KB \\
Records included & 6{,}825 \\
\midrule
\texttt{manifest\_sha\_match} (vs.\ \texttt{manifest.sha256}) & \checkmark True \\
\texttt{audit\_records.jsonl} SHA-256 match & \checkmark True \\
\texttt{chain\_verification.json} SHA-256 match & \checkmark True \\
\texttt{README.txt} SHA-256 match & \checkmark True \\
Hash chain integrity within bundle & \checkmark Valid \\
\midrule
\textit{Overall verification} & \checkmark Pass \\
\bottomrule
\end{tabular}
\end{table}

The bundle is reviewer-portable: a regulator with only \texttt{sha256sum} and a JSON parser can verify the manifest, every contents file, and the chain integrity. With \texttt{mr-audit} installed, \texttt{mr-audit inspect <bundle>} produces the same verification in one command.

We additionally verified that tampering with any record inside the bundle is detected. Modifying a single byte of \texttt{audit\_records.jsonl} (e.g., changing one prediction value from 0.142385 to 0.999) breaks the per-file SHA-256 against the manifest, and \texttt{verify\_bundle} returns \texttt{overall\_valid = False} with an explanatory error pointing at the changed file. The hash-chain integrity check inside \texttt{chain\_verification.json} also fails, providing a second line of defence.

\subsection{Summary of the demonstration}\label{sec:demonstration_summary}

The four use cases together demonstrate that a 6,825-record audit trail produced by \texttt{mr-audit} on a real seven-model volatility-forecasting pipeline can be:

\begin{itemize}[leftmargin=*]
    \item Replayed exactly (use case 1, 5/5 sampled records match bit-identically),
    \item Monitored for distributional drift across calendar windows (use case 2, 44/44 months flagged consistently with the documented regime shift),
    \item Used as the sole basis for a per-model validation report (use case 3, 7 models, 975 predictions each),
    \item Exported as a regulator-portable, integrity-tied bundle (use case 4, 841 KB, all checks pass).
\end{itemize}

The demonstration's purpose is not to argue that the seven-model panel itself satisfies SR 11-7 (that determination is the bank's, not the package's); it is to show that the technical record-keeping that an SR 11-7-compliant deployment would require is producible and verifiable using vendor-neutral, open-source tooling at a realistic data scale.
